% !TEX root = ../math.tex
\chapter{Mastering Dotter}
\label{chap:practice}

\section{The Cost of Incorrect Signals}

Expertise in Dotter is about accuracy in three ways:
(1) The user must synchronize their signals with the target timers.
(2) The user must accurately select the correct prefix.
(3) The user must not make involuntary signals (e.g. involuntary blinks).

From the program's perspective, there is virtually no difference between the user selecting
the wrong prefix and making an involuntary signal. Both are considered as incorrect signals and both are assumed
to have random phase. If the user intended the correct timer and was merely imprecise in synchronizing their signal,
we do not consider this to be an ``incorrect'' signal.
\footnote{
    We may safely assume that an involuntary signal has random phase.
    As for an incorrectly selected prefix, if we assume (perhaps unreasonably) that
    the incorrect prefix was randomly selected from all prefixes, and we further assume (legitimately) that we 
    have uniformly distributed the prefixes' timers, then such a signal has random phase as well.
}.

Here we will quantify the cost of incorrect signals.

UNFINISHED
% TODO: properly write this as I(D; h)

The user's speed can be calculated from their performance parameters.
Here we consider the effect of $\rho$, the probability of
unintended signals.
Given an observation, $D$, the gain in logits for the true hypothesis
$h$ is $\ln P(D | h) - \ln P(D | \neg h)$. When the posterior probability of
$h$ is small, this gain is equal to the gain in the log posterior probability
of $h$.
Let $r$ denote the hypothesis that the signal was unintended, so that
$P(r) = \rho$. 

We state the following facts:
\begin{align*}
P(D | \neg h, r) = P(D | \neg h, \neg r) &= P(D | \neg h) \\
P(D | \neg h) &\approx P(D), \quad \text{when } P(h) \approx 0 \\
P(D | h, r) &= P(D),  \quad \text{assuming that $D$, $h$ are independent conditioned on $r$} \\
\end{align*}

The likelihood under the true hypothesis is
$$P(D | h) = \rho P(D | h, r) + (1 - \rho) P(D | h, \neg r)$$

We consider the case where $P > > \sigma$
so that the components of the mixture are well-separated.
This means that the contribution of the false component is negligible
and that we can approximate the likelihood $P(D|h)$ piecewise as follows:

\[
P(D|h) \approx 
\begin{cases}
P(D | h, r) P(r) & \text{if } D \sim P(D|h, r) \\
P(D | h, \neg r)P(\neg r) & \text{if } D \sim P(D|h, \neg r).
\end{cases}
\]

That is, when $D$ is likely drawn from the unintended distribution, $P(D|h) \approx P(D|h, r) P(r)$. When $D$ is likely drawn from the intended distribution, $P(D|h) \approx P(D|h, \neg r)P(\neg r)$.




\begin{align*}
E_{D | r, h} \left[ \ln P(D | h) \right] &\approx E_{D | r, h} \left[ \ln \Big( P(D | h, r) P(r) \Big) \right] \\
&= \ln \rho - h(D | h, r) \\
\end{align*}
and similarly
\begin{align*}
E_{D | \neg r, h} \left[ \ln P(D | h) \right] &\approx \ln (1 - \rho) - h(D | h, \neg r)
\end{align*}
The expected log likelihood of the data under the true hypothesis is
\begin{align*}
E_D \left[ \ln P(D | h) \right] &= \rho E_{D | r, h} \left[ \ln P(D | h) \right] + (1 - \rho) E_{D | \neg r, h} \left[ \ln P(D | h) \right] \\
&\approx \rho (\ln \rho - h(D | h, r)) + (1 - \rho) (\ln (1 - \rho) - h(D | h, \neg r)) \\
&= - H_2(\rho) - \rho h(D | h, r) - (1 - \rho) h(D | h, \neg r) \\
&= - H_2(\rho) - \rho h(D) - (1 - \rho) h(D | h, \neg r), \quad \text{assuming $D|r$, $h|r$, independent}
\end{align*}
Where $H_2$ is the binary entropy function.

The likelihood of the alternative hypothesis is
% separating cases here enlightens nothing
\begin{align*}
E_{D | h}\left[ \ln P(D | \neg h) \right] &\approx E_{D | h} \left[ \ln P(D) \right] \\
\end{align*}
To simplify this term we use its expectation over $h$.
Recalling the fact that,
$$ E_a \left[ E_{b | a} [f(b)] \right] = E_b [f(b)] $$,
we find that this term's expectation over $h$ is $-h(D)$.

So the total expected gain in the logits of our hypothesis is
\begin{align*}
E_{D | h}[\ln P(D | h) - P(D | \neg h)] &\approx - H_2(\rho) + (1 - \rho) h(D) - (1 - \rho) h(D | h, \neg r) \\
\end{align*}
In our particular model, $D$ is uniformly distributed over an interval of length $P$ and
the entropy of $D$ conditional on $h$ is just the entropy of the gaussian noise.
Denoting the latter term by $H_{\sigma}$, we have,

\begin{align*}
E_{D | h}[\ln P(D | h) - P(D | \neg h)] &\approx - H_2(\rho) + (1 - \rho) (\ln P - H_{\sigma})
\end{align*}

This has a nice interpretation: we gain $\ln P - H_{\sigma}$ bits of information
per intended signal, but the presence of unintended signals makes us pay a penalty in two ways:
only $(1 - \rho)$ of the time is our signal intended, and additionally we must
expend $H_2(\rho)$ bits to learn whether the signal was intended.

Unfortunately the singularity of the binary entropy function
at $\rho = 0$ becomes a practical problem for fast text entry. For example, a
seemingly low outlier rate of $\rho = 0.03$ imposes the relatively
high cost of $H_2(0.03) \approx 0.19$ bits in addition to the
$3\%$ loss that would be naively expected.
Under reasonable parameters for an intermediate user (e.g. $\sigma = 25ms$, $P = 700ms$),
these combined terms amount to a loss of more than $10\%$ of the channel's capacity.
Thus to achieve expert entry speeds, it is necessary to
make unintended signals very rarely.

\section{Zen of Dotter}
Mastery has begun to be attained when you feel that the circles are turning slowly,
and that you are waiting, rather than reacting.
Practicing at extreme speeds (e.g. 340ms period) can help to achieve this state at a more
relaxed tempo (e.g. 700ms period). I have observed this effect on myself to complete satisfaction
multiple times. 
You are looking ahead, and are not thinking about anything other than the phrase.
You feel yourself to be willing the phrase onto the screen and have ceased to 
be conscious of the blinks.

